\chapter{Composite Systems, Entanglement, and the Singlet}

\lectureoverview{A question about photon emission exposes the difference between a quantum event and its ensemble average. That distinction leads naturally to measurement, collapse, and the need to treat a system together with its apparatus. We therefore construct the tensor-product state space of two systems, specialize to two spins, distinguish product states from entangled states, and embed each one-spin observable into the composite space. The singlet state then provides the decisive example: each spin is locally unpolarized, yet equal components are perfectly anticorrelated. Its basis-independent form prepares the ground for Bell's theorem and for the distinction between nonclassical correlation and faster-than-light signaling.}

\section{Discrete Transitions from a Superposed Spin}

Begin with one spin in a magnetic field along the \(z\)-axis. In units with \(\hbar=1\),
\begin{equation}
H=\frac{\omega}{2}\sigma_z,
\qquad
E_+=+\frac{\omega}{2},
\qquad
E_-=-\frac{\omega}{2}.
\label{eq:l6-spin-hamiltonian}
\end{equation}
The two levels are separated by
\begin{equation}
\Delta E=E_+-E_-=\omega.
\label{eq:l6-energy-gap}
\end{equation}
When the radiation field is included, a spin in the upper level may fall to the lower level and emit a photon of energy \(\omega\). A spin already in the lower level has no lower spin state into which it can radiate.

\begin{figure}[htbp]
\centering
\includegraphics[width=0.86\textwidth]{lecture_06_figure_01.png}
\caption{The two spin-energy levels and their separation \(\omega\), with \(\hbar=1\).}
\lectureframe{00:02:35}
\end{figure}

The subtle case is a spin prepared sideways. A classical magnetic rotor has a continuously varying energy proportional to the cosine of its angle with the field. It precesses at angular frequency \(\omega\), radiates at that same frequency, and gradually falls toward its minimum-energy orientation. The total radiated energy depends continuously on the initial angle: an initially horizontal rotor releases half as much energy as one that starts at the top.

Quantum mechanics preserves that classical answer only as an average. Write
\begin{equation}
\ket{\psi}=\alpha\ket{u}+\beta\ket{d},
\qquad
|\alpha|^2+|\beta|^2=1,
\label{eq:l6-superposed-spin}
\end{equation}
where \(\ket{u}\) and \(\ket{d}\) are the upper and lower energy eigenstates. The upper component can decay; the lower component cannot. Thus
\begin{equation}
P(\text{one photon})=|\alpha|^2,
\qquad
P(\text{no photon})=|\beta|^2.
\end{equation}
If a photon is emitted, its energy is the full level spacing \(\omega\), not a fraction of it. The mean emitted energy is nevertheless
\begin{equation}
\expect{E_{\rm emitted}}=|\alpha|^2\omega.
\end{equation}
For a spin along \(x\), \(|\alpha|^2=|\beta|^2=1/2\), so the mean is \(\omega/2\). Each individual event still produces either one photon of energy \(\omega\) or no photon. The continuous classical result belongs to the ensemble average; discreteness belongs to each event.

\begin{classroomqa}
\audiencequestion{How precisely must an absorbed photon match the level spacing?}
\lecturerresponse{The resonance has a finite linewidth. If the excited state has lifetime \(\tau\), its characteristic energy width is of order \(\Delta E\sim\hbar/\tau\). A larger magnetic dipole moment generally couples more strongly to the radiation field, shortens the lifetime, and broadens the line. This is the energy-time uncertainty relation in the form relevant to unstable levels.}
\end{classroomqa}

\begin{classroomqa}
\audiencequestion{Can photon emission itself serve as an apparatus that measures spin?}
\lecturerresponse{It can reveal whether a decay occurred, but it is not the ideal repeatable measurement model needed here. A measurement of a spin component should leave the spin in the eigenstate associated with the recorded value. Decay instead drives every emitted branch to the lower level, so it does not simply record both possible input values while preserving their corresponding eigenstates.}
\end{classroomqa}

\section{Collapse and the Need for a Composite System}

Between interventions an isolated state evolves by the Schr\"odinger equation. Measurement is ordinarily stated as an additional rule. Let \(L\) have orthonormal eigenvectors \(\ket{\lambda_i}\), and expand
\begin{equation}
\ket{\psi}=\sum_i\alpha_i\ket{\lambda_i},
\qquad
L\ket{\lambda_i}=\lambda_i\ket{\lambda_i}.
\label{eq:l6-measurement-expansion}
\end{equation}
A measurement yields one value \(\lambda_k\) with probability \(|\alpha_k|^2\), after which the state is
\begin{equation}
\ket{\psi}\longrightarrow\ket{\lambda_k}.
\label{eq:l6-collapse}
\end{equation}
The disappearance of the other components from the post-measurement description is called collapse.

\begin{figure}[htbp]
\centering
\includegraphics[width=0.84\textwidth]{lecture_06_figure_02.png}
\caption{A state expanded in the eigenbasis of the observable \(L\), before one outcome is selected.}
\lectureframe{00:12:35}
\end{figure}

\begin{classroomqa}
\audiencequestion{What is wrong with the collapse rule?}
\lecturerresponse{Nothing is wrong with it as an operational rule for predicting measurement outcomes. The conceptual tension is that an apparatus is itself a physical system. If both object and apparatus obey quantum mechanics, why should their interaction require a second dynamical law in addition to unitary Schr\"odinger evolution? The route taken here is to describe object and apparatus as one larger quantum system.}
\end{classroomqa}

The sharp system-apparatus boundary was a central point of historical debate. Rather than settle that debate by declaration, we first learn how quantum systems are combined. Once the combined state space is available, measurement can be studied as an interaction that correlates the system with an apparatus. The immediate mathematical question is therefore:
\begin{quote}
Given two quantum systems, what is the state space of the whole?
\end{quote}

\section{The Tensor-Product State Space}

Call the systems \(A\) and \(I\), with Hilbert spaces \(\mathcal S_A\) and \(\mathcal S_I\). Their composite space is
\begin{equation}
\mathcal S_{AI}=\mathcal S_A\otimes\mathcal S_I.
\label{eq:l6-tensor-space}
\end{equation}
Choose orthonormal bases \(\{\ket{a}\}\) and \(\{\ket{i}\}\) for the two factors. Pair every basis vector of the first system with every basis vector of the second:
\begin{equation}
\ket{a}\otimes\ket{i}\equiv\ket{a\,i}.
\end{equation}
If the factor spaces have dimensions \(D_A\) and \(D_I\), then
\begin{equation}
\dim\mathcal S_{AI}=D_A D_I.
\end{equation}
The inner product factorizes,
\begin{equation}
\braket{b\,j}{a\,i}
=\braket{b}{a}\braket{j}{i}
=\delta_{ba}\delta_{ji}.
\label{eq:l6-tensor-inner-product}
\end{equation}
Thus two composite basis states are orthogonal whenever either subsystem label differs.

\begin{figure}[htbp]
\centering
\includegraphics[width=0.86\textwidth]{lecture_06_figure_03.png}
\caption{Separate state spaces, their tensor product, and the paired orthonormal basis.}
\lectureframe{00:24:05}
\end{figure}

The paired basis vectors are not the whole space. Their complete linear span is:
\begin{equation}
\ket{\Psi}=\sum_{a,i}c_{ai}\ket{a\,i},
\qquad
\sum_{a,i}|c_{ai}|^2=1.
\label{eq:l6-general-tensor-state}
\end{equation}

\begin{classroomqa}
\audiencequestion{Are \(\ket{a}\) and \(\ket{i}\) arbitrary vectors?}
\lecturerresponse{Here they are members of chosen orthonormal bases. Any orthonormal basis may be used in either factor. Pairing those basis vectors constructs an orthonormal basis of the composite space, after which arbitrary vectors are obtained by taking linear combinations as in Equation~\eqref{eq:l6-general-tensor-state}.}
\end{classroomqa}

\begin{classroomqa}
\audiencequestion{Does the tensor product depend on the chosen bases?}
\lecturerresponse{Its coordinate description does, but the space does not. Changing basis in either subsystem merely changes the induced basis used to describe the same vectors in \(\mathcal S_A\otimes\mathcal S_I\).}
\end{classroomqa}

\section{Two Spins and Product States}

Now take two spins. Denote the components of the first by \(\sigma_x,\sigma_y,\sigma_z\), and those of the second by \(\tau_x,\tau_y,\tau_z\). In a \(z\)-basis, the composite basis is
\begin{equation}
\ket{uu},\qquad \ket{ud},\qquad \ket{du},\qquad \ket{dd}.
\label{eq:l6-two-spin-basis}
\end{equation}
The first slot always refers to \(\sigma\), the second to \(\tau\). A general state is
\begin{equation}
\ket{\Psi}
=c_{uu}\ket{uu}+c_{ud}\ket{ud}
+c_{du}\ket{du}+c_{dd}\ket{dd}.
\label{eq:l6-general-two-spin}
\end{equation}

\begin{classroomqa}
\audiencequestion{If the two spins interact, do the basis vectors themselves change?}
\lecturerresponse{No fixed basis is required to follow the dynamics. In the chosen basis, interaction changes the coefficients \(c_{ab}(t)\), not the definition of \(\ket{uu},\ket{ud},\ket{du},\ket{dd}\). A state that begins as \(\ket{uu}\) can evolve into a superposition, or even into \(\ket{dd}\), while the coordinate basis remains the same.}
\end{classroomqa}

The construction must also pass a consistency test. If the spins are far apart and only \(\sigma\) is examined, the predictions must agree with the one-spin theory developed before \(\tau\) was introduced. Tensor products are designed to retain that local theory.

Some two-spin states describe independently prepared spins. For example, let
\begin{equation}
\ket{R}=\frac{\ket{u}+\ket{d}}{\sqrt2}
\end{equation}
be the \(+1\) eigenstate of the \(x\)-component. Then
\begin{equation}
\ket{u}_{\sigma}\otimes\ket{R}_{\tau}
=\frac{\ket{uu}+\ket{ud}}{\sqrt2}.
\label{eq:l6-up-right-product}
\end{equation}

\begin{classroomqa}
\audiencequestion{May the two factors be prepared along different axes?}
\lecturerresponse{Yes. Each factor can be any one-spin state. Equation~\eqref{eq:l6-up-right-product} has the first spin definite along \(z\) and the second definite along \(x\). Expanding the second factor in the \(z\)-basis merely rewrites the same product state in the composite basis.}
\end{classroomqa}

The most general product state is
\begin{align}
\ket{\Psi_{\rm prod}}
&=(\alpha_u\ket{u}+\alpha_d\ket{d})
\otimes(\beta_u\ket{u}+\beta_d\ket{d})\\
&=\alpha_u\beta_u\ket{uu}
+\alpha_u\beta_d\ket{ud}
+\alpha_d\beta_u\ket{du}
+\alpha_d\beta_d\ket{dd}.
\label{eq:l6-product-expansion}
\end{align}

\begin{figure}[htbp]
\centering
\includegraphics[width=0.86\textwidth]{lecture_06_figure_04.png}
\caption{A factorized two-spin state expanded in the four-state composite basis.}
\lectureframe{00:37:45}
\end{figure}

\begin{classroomqa}
\audiencequestion{What kind of vector product appears in Equation~\eqref{eq:l6-product-expansion}?}
\lecturerresponse{The ambient space is the tensor product of the two Hilbert spaces. A vector that itself factors into one vector from each space is called a product state or simple tensor. Most vectors in a tensor-product space are sums of such products and cannot be written as one product.}
\end{classroomqa}

The restriction can be seen directly in the coefficients:
\begin{equation}
c_{uu}c_{dd}=c_{ud}c_{du}
\qquad\text{for every product state}.
\label{eq:l6-rank-one-condition}
\end{equation}
Equivalently, the \(2\times2\) coefficient matrix has rank one. Parameter counting reaches the same conclusion. Four complex coefficients contain eight real parameters; normalization removes one and the irrelevant overall phase removes another, leaving six for a general pure two-spin state. Each pure one-spin state needs two real parameters, so two separately specified spins need only four. Product states therefore occupy a special four-parameter subset of the six-parameter two-spin state space. The remaining states are entangled.

\section{Local Observables in the Composite Space}

The familiar one-spin observables must survive inside the larger system. More explicitly,
\begin{equation}
\sigma_m\equiv\sigma_m\otimes I,
\qquad
\tau_n\equiv I\otimes\sigma_n,
\qquad
m,n\in\{x,y,z\}.
\label{eq:l6-local-operators}
\end{equation}
The abbreviated notation is harmless once the slot convention is fixed. Using
\begin{equation}
\sigma_x=
\begin{pmatrix}0&1\\1&0\end{pmatrix},
\qquad
\sigma_y=
\begin{pmatrix}0&-i\\i&0\end{pmatrix},
\qquad
\sigma_z=
\begin{pmatrix}1&0\\0&-1\end{pmatrix},
\end{equation}
the first-spin operators act while the second label is a spectator:
\begin{align}
\sigma_x\ket{u\,i}&=\ket{d\,i},
&
\sigma_x\ket{d\,i}&=\ket{u\,i},\\
\sigma_z\ket{u\,i}&=\ket{u\,i},
&
\sigma_z\ket{d\,i}&=-\ket{d\,i}.
\end{align}
For the second spin the first label is the spectator:
\begin{align}
\tau_x\ket{a\,u}&=\ket{a\,d},
&
\tau_x\ket{a\,d}&=\ket{a\,u},\\
\tau_z\ket{a\,u}&=\ket{a\,u},
&
\tau_z\ket{a\,d}&=-\ket{a\,d}.
\label{eq:l6-spectator-action}
\end{align}
The \(y\)-component also flips its active slot, with
\(\sigma_y\ket{u}=i\ket{d}\) and
\(\sigma_y\ket{d}=-i\ket{u}\).

\begin{figure}[htbp]
\centering
\includegraphics[width=0.86\textwidth]{lecture_06_figure_05.png}
\caption{The second-spin action table: the first index remains a spectator.}
\lectureframe{00:47:40}
\end{figure}

\begin{classroomqa}
\audiencequestion{Does applying \(\sigma_x\) mean that an \(x\)-measurement has occurred, and why is the result not a superposition?}
\lecturerresponse{Operator action is not itself a measurement. Multiplying the matrix \(\sigma_x\) by the \(z\)-basis vector \(\ket{u}\) gives \(\ket{d}\). This is not an eigenvalue equation: an eigenvector would keep its direction under the action. Indeed, \(\ket{u}\) is mapped to an orthogonal state, which is another way to see why an \(x\)-measurement on a \(z\)-eigenstate is maximally random. Acting on a non-eigenvector is usually an intermediate algebraic step whose physical meaning comes from the full quantity being calculated.}
\end{classroomqa}

Operators on different factors commute. For example,
\begin{equation}
\sigma_x\tau_x\ket{uu}
=\ket{dd}
=\tau_x\sigma_x\ket{uu}.
\end{equation}
The same equality holds on every composite basis vector and therefore on every state:
\begin{equation}
[\sigma_m,\tau_n]=0
\qquad\text{for all }m,n.
\label{eq:l6-local-commutation}
\end{equation}

\begin{figure}[htbp]
\centering
\includegraphics[width=0.86\textwidth]{lecture_06_figure_06.png}
\caption{Separate slot actions lead to the vanishing commutator \([\sigma_x,\tau_x]\).}
\lectureframe{00:56:40}
\end{figure}

\begin{classroomqa}
\audiencequestion{Why is checking one basis vector not yet a proof that two operators commute?}
\lecturerresponse{Operator equality requires equal action on every vector. It is enough to check every vector of a basis, because linearity then extends the result to every superposition. For local operators the argument is immediate on all four basis states: each operator changes a different slot, so their order is irrelevant.}
\end{classroomqa}

\begin{classroomqa}
\audiencequestion{Are there observables that belong genuinely to the two-spin system?}
\lecturerresponse{Yes. Products such as \(\sigma_x\tau_x\), \(\sigma_x\tau_y\), and linear combinations of such products probe joint properties. The local observables describe each factor; correlation observables reveal how the factors are related.}
\end{classroomqa}

\begin{classroomqa}
\audiencequestion{Are independently prepared product states already entangled?}
\lecturerresponse{No. In a product state each factor has its own state vector, and local expectation values factorize. Entanglement begins when the composite state cannot be separated into one state for \(\sigma\) and one for \(\tau\).}
\end{classroomqa}

\section{The Singlet State}

A maximally entangled example is
\begin{equation}
\ket{S}
=\frac{1}{\sqrt2}\bigl(\ket{ud}-\ket{du}\bigr).
\label{eq:l6-singlet}
\end{equation}
The two terms already show perfect \(z\)-anticorrelation: whenever one spin is up, the other is down. The deeper point is that neither spin can be assigned its own pure state.

\begin{figure}[htbp]
\centering
\includegraphics[width=0.86\textwidth]{lecture_06_figure_07.png}
\caption{The antisymmetric two-spin state that will be identified as the singlet.}
\lectureframe{01:09:25}
\end{figure}

For a pure state of one spin, there is always a unit vector \(\bm n\) for which
\begin{equation}
(\bm\sigma\cdot\bm n)\ket{\psi}=\ket{\psi}.
\end{equation}
Equivalently, the Bloch vector
\begin{equation}
\bm r
=\bigl(\expect{\sigma_x},\expect{\sigma_y},\expect{\sigma_z}\bigr)
\end{equation}
has length one for a pure one-spin state. Its three components therefore cannot all vanish. Given the state, orienting the detector along \(\bm r\) yields \(+1\) with certainty. The familiar examples are \(z\) for \(\ket{u}\), \(x\) for \((\ket{u}+\ket{d})/\sqrt2\), and one of the \(y\)-directions for \((\ket{u}+i\ket{d})/\sqrt2\).

Now calculate a local average in the singlet:
\begin{align}
\expect{\sigma_z}_S
&=\frac12
(\bra{ud}-\bra{du})\sigma_z
(\ket{ud}-\ket{du})\\
&=\frac12
(\bra{ud}-\bra{du})
(\ket{ud}+\ket{du})\\
&=0.
\end{align}
For \(\sigma_x\),
\begin{equation}
\sigma_x\ket{S}
=\frac{\ket{dd}-\ket{uu}}{\sqrt2},
\end{equation}
which is orthogonal to \(\ket{S}\), so \(\expect{\sigma_x}_S=0\). The phases supplied by \(\sigma_y\) do not change that orthogonality, hence
\begin{equation}
\expect{\sigma_x}_S
=\expect{\sigma_y}_S
=\expect{\sigma_z}_S=0.
\label{eq:l6-zero-local-sigma}
\end{equation}
By symmetry,
\begin{equation}
\expect{\tau_x}_S
=\expect{\tau_y}_S
=\expect{\tau_z}_S=0.
\label{eq:l6-zero-local-tau}
\end{equation}

\begin{figure}[htbp]
\centering
\includegraphics[width=0.86\textwidth]{lecture_06_figure_08.png}
\caption{The board records vanishing local spin averages and the \(y\)-component flip used to complete the check.}
\lectureframe{01:14:20}
\end{figure}

The whole two-spin state is pure and perfectly known, but each part is locally unpolarized. Since no pure one-spin state has a zero Bloch vector, \(\ket{S}\) cannot be a product state.

\begin{classroomqa}
\audiencequestion{Would an ordinary product state still have a direction with nonzero local polarization?}
\lecturerresponse{Yes. The one-spin proof applies to the relevant factor, while the other factor is only a spectator. Thus each factor of a pure product state has a unit Bloch vector. The singlet fails this product-state property for both spins.}
\end{classroomqa}

\begin{classroomqa}
\audiencequestion{How can two spins be prepared in the singlet?}
\lecturerresponse{A coupling that favors anti-alignment can make the singlet the lowest-energy state. Bring the spins into interaction and allow the pair to lose energy to its surroundings; it can relax into the singlet. The statement depends on the interaction. For the isotropic antiferromagnetic Hamiltonian \(H_{\rm int}=J\,\bm\sigma\cdot\bm\tau\) with \(J>0\), the singlet lies below the three triplet states.}
\end{classroomqa}

\begin{classroomqa}
\audiencequestion{How long does that relaxation take?}
\lecturerresponse{The timescale depends on the separation of the spins, the strength of their mutual coupling, their magnetic moments, and their coupling to the electromagnetic environment that carries away energy. Stronger coupling generally means faster relaxation.}
\end{classroomqa}

\begin{classroomqa}
\audiencequestion{Why can the antisymmetric combination have lower energy than the symmetric one?}
\lecturerresponse{For two spin-\(\tfrac12\) particles, the antisymmetric state has total spin zero, while the three symmetric states form a total-spin-one triplet. An interaction proportional to \(+\bm\sigma\cdot\bm\tau\) assigns eigenvalue \(-3\) to the singlet and \(+1\) to the triplet, so positive coupling places the singlet lower. Other Hamiltonians need not order the levels the same way.}
\end{classroomqa}

\subsection{The Same State in Every Spin Basis}

Entanglement is not removed by choosing a more convenient local basis. A basis-independent measure, developed later, is the entanglement entropy. The singlet displays an even stronger symmetry. With
\begin{equation}
\ket{u}=\frac{\ket{R}+\ket{L}}{\sqrt2},
\qquad
\ket{d}=\frac{\ket{R}-\ket{L}}{\sqrt2},
\end{equation}
substitution into Equation~\eqref{eq:l6-singlet} gives
\begin{equation}
\ket{S}
=\frac{\ket{LR}-\ket{RL}}{\sqrt2}.
\label{eq:l6-singlet-x-basis}
\end{equation}
In any common orthonormal spin basis \(\{\ket{\rm in},\ket{\rm out}\}\),
\begin{equation}
\ket{S}
=\frac{\ket{{\rm in}\,{\rm out}}-\ket{{\rm out}\,{\rm in}}}{\sqrt2},
\label{eq:l6-singlet-any-basis}
\end{equation}
up to an irrelevant overall phase. The three triplet states mix among themselves under a common rotation; the one-dimensional singlet sector can only reproduce itself.

\begin{figure}[htbp]
\centering
\includegraphics[width=0.86\textwidth]{lecture_06_figure_09.png}
\caption{The antisymmetric pattern survives changes from up-down to right-left and other common spin bases.}
\lectureframe{01:21:00}
\end{figure}

\begin{classroomqa}
\audiencequestion{Is entanglement merely an artifact of the basis used to write the state?}
\lecturerresponse{No. Local basis changes alter the coefficients but not whether the state factorizes, and entanglement entropy quantifies that invariant fact. For the singlet, common rotations make the invariance especially visible: the antisymmetric form is reproduced in every spin basis.}
\end{classroomqa}

\section{Correlation Observables}

Local averages vanish, so the information in the singlet must reside in joint observables. Start with
\begin{equation}
\sigma_z\tau_z\ket{ud}=-\ket{ud},
\qquad
\sigma_z\tau_z\ket{du}=-\ket{du}.
\end{equation}
Therefore
\begin{equation}
\sigma_z\tau_z\ket{S}=-\ket{S},
\qquad
\expect{\sigma_z\tau_z}_S=-1.
\label{eq:l6-zz-correlation}
\end{equation}
The result is not merely an average: the singlet is an eigenstate of the product, so equal-axis \(z\)-measurements always have opposite signs.

Both \(x\)-operators flip their slots:
\begin{equation}
\sigma_x\tau_x\ket{ud}=\ket{du},
\qquad
\sigma_x\tau_x\ket{du}=\ket{ud}.
\end{equation}
The exchange of the two terms reverses the antisymmetric state:
\begin{equation}
\sigma_x\tau_x\ket{S}=-\ket{S},
\qquad
\expect{\sigma_x\tau_x}_S=-1.
\end{equation}
The two phases from the \(y\)-operators cancel in the same way, giving
\begin{equation}
\expect{\sigma_y\tau_y}_S=-1.
\end{equation}
By contrast, unequal Cartesian components have zero expectation. For example,
\begin{equation}
\sigma_x\tau_z\ket{S}
=-\frac{\ket{dd}+\ket{uu}}{\sqrt2},
\qquad
\expect{\sigma_x\tau_z}_S=0.
\end{equation}
All these results are summarized by
\begin{equation}
\expect{\sigma_m\tau_n}_S=-\delta_{mn},
\qquad
\expect{(\bm\sigma\cdot\bm a)(\bm\tau\cdot\bm b)}_S
=-\bm a\cdot\bm b.
\label{eq:l6-singlet-correlation-tensor}
\end{equation}

\begin{figure}[htbp]
\centering
\includegraphics[width=0.86\textwidth]{lecture_06_figure_10.png}
\caption{The correlation calculation records \(\expect{\sigma_x\tau_x}=-1\) and the mixed-axis result \(\expect{\sigma_x\tau_z}=0\).}
\lectureframe{01:32:20}
\end{figure}

\begin{classroomqa}
\audiencequestion{What does a product of two observables mean physically? Is one measurement performed after the other?}
\lecturerresponse{In general a Hermitian product is another observable on the composite space. Here every \(\sigma_m\) commutes with every \(\tau_n\), so the two local components may be measured simultaneously and their numerical outcomes multiplied. The product operator packages that joint result; no time ordering is required.}
\end{classroomqa}

For a product state, joint averages factorize:
\begin{equation}
\expect{\sigma_m\tau_n}_{\rm prod}
=\expect{\sigma_m}_{\sigma}\expect{\tau_n}_{\tau}.
\end{equation}
The singlet violates this factorization maximally. Each factor on the right would be zero, while an equal-axis product on the left is \(-1\). Individually the spins are completely random; jointly, equal components are certainly opposite.

\begin{classroomqa}
\audiencequestion{Does every interaction between two spins produce entanglement?}
\lecturerresponse{A generic coupling can entangle an initially factorized state, and many ordinary interactions do so. It is not literally inevitable: a specially chosen product eigenstate may only acquire a phase, and a purely local Hamiltonian cannot entangle the factors. What matters is a nontrivial coupling acting on a state that it does not preserve as a product.}
\end{classroomqa}

\section{From Correlation to Bell's Question}

Perfect anticorrelation alone is not paradoxical. Imagine placing a penny and a dime in a pocket, distributing one unseen coin to Alice and the other to Bob, and sending them far apart. When Alice sees which coin she has, she immediately knows Bob's coin. The correlation was established when the coins were together; learning one result sends no signal to the other location.

The singlet shares that elementary feature but adds a crucial freedom: Alice and Bob may choose any measurement directions, and quantum mechanics predicts the continuously varying correlation \(-\bm a\cdot\bm b\). Bell's theorem asks whether all those correlations can arise from local classical instructions carried by the particles in advance.

\begin{classroomqa}
\audiencequestion{Is this the state used in Bell's theorem?}
\lecturerresponse{Yes. Bell-type arguments use entangled pairs such as the singlet, not a definite product state \(\ket{ud}\). The distinction matters because a product state supplies separate local descriptions, whereas the singlet carries one inseparable correlation structure across every common basis.}
\end{classroomqa}

\begin{classroomqa}
\audiencequestion{Why call the anticorrelation statistical if equal-axis outcomes are always opposite?}
\lecturerresponse{Each local outcome is random: Alice cannot predict whether she will obtain \(+1\) or \(-1\). The conditional relation is exact: once one equal-axis outcome is known, the other is fixed to be opposite. Correlation refers to the joint distribution, not to uncertainty in the equal-axis product.}
\end{classroomqa}

A useful way to sharpen the issue is through simulation. One classical computer can imitate every experiment on one spin: store its state, evolve it with the Schr\"odinger equation, calculate the requested probabilities, and use a random-number generator to select an outcome. Repeated measurements and magnetic precession can all be reproduced.

Now split an entangled pair between two classical computers. Allow the computers to interact while the pair is prepared, then disconnect them and choose measurement directions independently at distant locations. No pair of separated local classical programs, even with shared randomness fixed in advance, can reproduce the quantum correlations for every choice of directions. A global classical simulation is possible, but it must keep access to the joint state or otherwise coordinate the two sides in a way that is not a local hidden-variable model.

This failure of local classical bookkeeping does not provide a faster-than-light communication channel. Alice's local outcomes remain random, so she cannot control a message that Bob could read without later comparing their records. The terminology surrounding nonlocality varies, but the operational distinction is firm: Bell correlations rule out local classical explanations while preserving no-signaling.

The composite-state construction has therefore produced something genuinely new. Tensor products preserve the old one-spin physics locally, yet their generic vectors cannot be divided into separate states. In the singlet, all local polarization vanishes while joint spin components carry perfect, basis-independent anticorrelation. The next step is to turn this structure into Bell's quantitative test of classical reasoning.
